\documentclass[../main.tex]{subfiles}

\begin{document}

% Eksempler på de mest brugte ting i skabelonen. Slet kapitlet (og linjen i main.tex), når du ikke har brug for det.

\section{Eksempler}
\label{sec:eksempler}

\subsection{Tekst, forkortelser og kilder}
\label{subsec:tekst}
Forkortelser skrives med \verb|\ac{au}|: første gang \ac{au}, derefter \ac{au}.
Kilder angives med \verb|\cite{AUWebsite}| og bliver til en fodnote\cite{AUWebsite}.
Henvisninger laves med \verb|\cref|: se \cref{fig:segl}, \cref{tab:eksempel}, \cref{eq:ohm}, \cref{lst:hej} og \cref{app:eksempel}.
Enheder skrives med siunitx: \qty{6,5}{\kilo\watt} og \num{1234,5}.
\todo{En note i margenen}

\subsection{Figur}
\label{subsec:figur}
\begin{figure}[H]
    \centering
    \includegraphics[width=0.3\linewidth]{auSegl.pdf} % billeder lægges i Billeder/ og skrives som {navn.png}
    \caption{Aarhus Universitets segl}
    \label{fig:segl}
\end{figure}

\subsection{Tabel}
\label{subsec:tabel}
\begin{table}[H]
    \centering
    \renewcommand{\arraystretch}{1.6} % højere rækker, som på forsiden
    \begin{tabular}{|l|r|r|}
        \hline
        \rowcolor{lightgray}
        \textbf{Komponent} & \textbf{Antal} & \textbf{Pris} \\
        \hline
        Modstand  & 10 & \qty{0,50}{kr} \\
        \hline
        Kondensator & 4 & \qty{1,25}{kr} \\
        \hline
    \end{tabular}
    \caption{Eksempel på en tabel}
    \label{tab:eksempel}
\end{table}

\subsection{Formel}
\label{subsec:formel}
\begin{equation}
    U = R \cdot I
    \label{eq:ohm}
\end{equation}

\subsection{Kode}
\label{subsec:kode}
\begin{lstlisting}[caption={Hello world i C++}, label={lst:hej}]
#include <iostream>

int main() {
    std::cout << "Hej verden" << std::endl; // udskriv
    return 0;
}
\end{lstlisting}

\end{document}
